\begin{table}[htbp]
    \centering
    \caption{Brechas operativas diagnosticadas en el proceso de inscripción del MTP San Julián (gestiones 2022 a 2025).}
    \label{tab:brechas_diagnostico}
    \begin{tabularx}{\textwidth}{L L L}
        \toprule
        \textbf{Componente} & \textbf{Estado diagnosticado} & \textbf{Evidencia (2025-2026)} \\
        \midrule
        Recuperación de credenciales & Sin autoservicio; restablecimiento manual presencial en ventanilla & Cientos de reseteos por gestión; pérdida de cupos vocacionales mientras el postulante permanece en fila \\
        Matriculación documental definitiva & Sin módulo informático; archivo de carpetas físicas sin base de datos & 0\% de registros auditables en 2025; expedientes incompletos y sobrecarga de secretaría académica \\
        Asignación a cursos y paralelos por unidad educativa & Distribución manual en planillas externas según colegio de procedencia & Errores de transcripción; sobrepoblación en talleres y conflictos de coordinación horaria \\
        \bottomrule
    \end{tabularx}
    \notatabla{Elaboración propia con base en el análisis retrospectivo de las gestiones 2022 a 2025 y el relevamiento de campo del primer bimestre de 2026 (60 estudiantes postulantes y 5 funcionarios) del MTP San Julián. Población de referencia: cerca de 400 postulantes anuales en siete especialidades.}
\end{table}
